\section{Results}
\label{sec:results}

This section presents the empirical results of both experiments across four phases of increasing model complexity. In Experiment~1, we vary the number of expert-labelled samples $n \in [20, 200]$ while keeping the total dataset size fixed at $N = 997$. In Experiment~2, we fix $n \in \{50, 100, 200\}$ and vary the total pool size $N$. All results are averaged over 300 random repetitions. Performance is measured by standardised RMSE (sRMSE) and standardised bias relative to $\theta^*$, the MLE fit on all expert labels.

The five annotating LLMs are \textsc{Llama-3.1-8B-Instruct}, \textsc{Mistral-7B-Instruct-v0.3}, \textsc{DeepSeek-V3}, \textsc{GPT-5.4}, and \textsc{Claude Sonnet 4.5}. The four datasets are CUAD (binary legal clause detection), Misogynistic (binary hate speech), VUAMC (binary metaphor detection), and FOMC (binary monetary policy stance).

%% ============================================================
\subsection{Experiment 1: Varying Expert Sample Size}
\label{sec:exp1}
%% ============================================================

\subsubsection{Phase 1 — Class Prevalence (Intercept-Only)}
\label{sec:exp1-phase1}

In the simplest setting — intercept-only logistic regression estimating class prevalence — all three debiasing methods consistently outperform the expert-only baseline across all datasets and LLMs. PPI and PPI++ achieve the lowest sRMSE overall, with DSL converging to a similar level around $n \approx 33$. Bias curves are jagged due to discrete sampling but converge toward zero for all methods as $n$ grows.

The benefit of debiasing is clearest on FOMC, where all methods beat the LLM baseline for every LLM, and on VUAMC and Misogynistic, where most methods beat the LLM baseline when the annotating LLM has lower agreement with expert labels. On CUAD, LLM annotation quality is so high that debiasing offers only marginal improvements. Across LLMs, the relative method ordering is stable, suggesting annotation quality has limited impact in this low-complexity setting.

\subsubsection{Phases 2 and 3 — Single Feature}
\label{sec:exp1-phase23}

Adding a single covariate (low-variance in Phase~2, high-variance in Phase~3) degrades PPI's relative performance more than the other methods: at small $n$, PPI frequently starts \emph{above} the expert-only baseline before recovering. DSL and PPI++ are more stable throughout, with PPI++ generally best at very low $n$ and DSL converging tightly as $n$ grows. The methods still beat the expert-only baseline on CUAD and FOMC, but on VUAMC no method consistently beats expert-only, the first sign that VUAMC poses a fundamentally harder debiasing problem.

Misogynistic shows a strong Phase~3 pattern: PPI never beats the expert-only baseline and has a large gap in both sRMSE and bias, while DSL and PPI++ overlap tightly and both beat the LLM baseline for all LLMs. This behaviour is consistent across all five LLMs, pointing to a dataset-specific rather than LLM-specific mechanism.

\subsubsection{Phase 4 — Full Logistic Regression}
\label{sec:exp1-phase4}

The full five-feature model is the most practically relevant setting. The broad patterns from Phases 2–3 persist, but certain phenomena become more pronounced:

\begin{itemize}
    \item \textbf{PPI++ is the most reliable method.} PPI++ beats the expert-only baseline on CUAD, Misogynistic, and FOMC across virtually all LLMs, and performs best or tied-best in bias.
    \item \textbf{PPI consistently underperforms at low $n$.} On every dataset, PPI starts at a much higher sRMSE than the other methods and only recovers as $n$ grows. On Misogynistic it never recovers to beat expert-only. This instability is most severe for LLMs with lower annotation quality.
    \item \textbf{DSL goes flat on CUAD and Misogynistic.} After approximately $n/N \approx 0.04$, DSL's sRMSE stops decreasing despite additional expert annotations, suggesting the regularizer dominates the LLM correction signal. This is not observed for PPI or PPI++.
    \item \textbf{VUAMC remains the hardest dataset.} No method beats the expert-only baseline consistently. PPI++ is the best debiasing method and beats expert-only at low $n$ (up to roughly $n/N \approx 0.13$), but the gain disappears as $n$ grows. On FOMC, DSL and PPI++ are both competitive with PPI++ leading on bias.
\end{itemize}

%% ============================================================
\subsection{Experiment 2: Varying Total Pool Size}
\label{sec:exp2}
%% ============================================================

Experiment~2 asks whether access to more unlabelled data (larger $N$) improves debiasing for a fixed expert budget $n$. The key finding is that the debiasing methods benefit from larger $N$, but the benefit is uneven across datasets and phases, and it diminishes as $n$ grows.

Across Phases 1–3, increasing $N$ generally lowers sRMSE for all debiasing methods, but expert-only sRMSE is unaffected (it depends only on $n$). This means the gap between debiasing methods and expert-only tends to grow with $N$ — a positive result. The LLM baseline (fit on the full unlabelled pool) decreases with $N$ as expected, and the debiasing methods beat the LLM baseline more consistently at larger $N$.

In Phase~4 (full logistic), the relative ordering of methods mirrors Experiment~1: PPI++ leads, DSL is stable and parallel to PPI++, and PPI starts highest and converges. On CUAD and FOMC, DSL and PPI++ beat the expert-only baseline across all tested $N$. On VUAMC, DSL and PPI++ are nearly flat as a function of $N$ and remain below the expert-only baseline for $n \geq 100$; at $n = 200$, PPI++ and expert-only become indistinguishable. On Misogynistic, PPI remains consistently above the expert-only baseline regardless of $N$, while DSL and PPI++ both beat it.

A notable pattern in both experiments is the split between open-source LLMs (Llama, Mistral, DeepSeek) and closed-source LLMs (GPT, Claude): closed-source models have uniformly lower LLM baselines on most datasets, meaning the debiasing correction is harder to beat. Despite this, the relative ordering of the debiasing methods themselves is stable across LLM families.

%% ============================================================
\subsection{General Observations}
\label{sec:general-obs}
%% ============================================================

Several patterns hold consistently across both experiments, all phases, and all datasets:

\begin{enumerate}
    \item \textbf{All debiasing methods outperform expert-only at small $n$,} except on VUAMC in higher phases. The advantage is largest at $n \approx 20$–50 and diminishes as $n$ grows — consistent with the theoretical motivation for semi-supervised correction.

    \item \textbf{PPI++ is the most robust method overall.} It is stable across all model complexities, achieves competitive sRMSE and bias, and never exhibits the large instabilities seen with PPI. Its advantage over DSL is most visible at very small $n$.

    \item \textbf{PPI is sensitive to annotation quality and model complexity.} It performs best in Phase~1 and deteriorates relative to the other methods as features are added. Its poor performance is most pronounced on datasets where LLM annotation quality is lower (Llama on VUAMC and FOMC), consistent with PPI's correction term amplifying rather than reducing error when $\hat{Y}$ is noisy.

    \item \textbf{DSL is stable but can plateau.} In the full logistic phase, DSL shows flat convergence on CUAD and Misogynistic, suggesting the regularizer dominates. Its bias performance is generally good.

    \item \textbf{VUAMC is the most challenging dataset.} It is the only dataset for which no method consistently beats expert-only across phases and experiments. A likely explanation is that metaphor detection relies on contextual cues that LLMs annotate inconsistently, so pseudo-labels add noise rather than signal.

    \item \textbf{The choice of LLM affects the scale of improvement but not the method ordering.} Higher annotation quality (e.g.\ Claude and GPT on CUAD) compresses the sRMSE scale, but the relative ranking of DSL, PPI, and PPI++ is preserved across LLMs within each dataset and phase.
\end{enumerate}

\paragraph{Open questions.} The results leave several patterns unexplained that warrant further investigation: why PPI degrades faster than the other methods with model complexity; why DSL plateaus in the full logistic phase on specific datasets; and whether the annotation quality metric (Cohen's $\kappa$) is a reliable predictor of which method benefits most. A formal correlation analysis between $\kappa$ and per-method sRMSE improvement, as well as a bias-variance decomposition, would help address these questions.
